\section{账本事件的区域史深描：mongol-yuan}

\label{sec:ledger-depth-mongol-yuan}

本组把账本事件置于战争、制度、区域社会与证据类型的交叉处。年表提供政治节点，地方材料与物质遗存则限制我们对征发、身份、信仰和迁徙所能作出的断言。

\subsection{战争、行政与区域资源}

账本条目\texttt{YUAN-1206-EVENT-151}所记1206年的“铁木真获成吉思汗称号，蒙古诸部统一”，属于蒙古的历史语境。条目所示前因是前期边防、财政和统治联盟的压力构成直接背景，具体条件须按本条材料分辨，其可见后果为改变后续资源配置与政治选择；实际执行范围和社会影响仍须分项核验。这一变化若进入区域生活，必须通过道路、文书、军队、市场、家庭和地方中介才会发生；政权易手或法令发布不等于所有城镇、乡村与边地在同一时间获得同样经验。

本条所列来源为《元史》〈本纪〉及相关志；《元朝秘史》；《元典章》或《通制条格》相关条。它们能够支持特定的年序、制度语言、政治行动或局部遗存，却不能无条件化为社会总量。账本特别提醒“编年主干可靠；细节、规模与动机须与对方材料、地方文书或考古材料互校”。因此，军报的兵数、条约的名分、官府的族群分类、判牍中的道德判断以及后出的叙述都须回到产生它们的场景，而不可直接代表全域动机或人口。

将事件与地方层次连接时，应同时考察资源怎样被征集、信息怎样传递、迁徙怎样改变户籍和劳动、宗教与亲属网络怎样提供或限制协作。现有证据往往只照亮少数地点；保留这种不均衡，才能避免把宋辽夏金元的多政权历史改写成单一中心的必然过程。

账本条目\texttt{YUAN-1206-REGNAL-YEAR}所记1206年的“蒙古以成吉思汗在位第1年作为诏令、赋役与外交文书的纪年框架”，属于蒙古的历史语境。条目所示前因是新岁颁历、年号和纪年是中枢与地方行政沟通的共同前提，其可见后果为为同年政令、战事和地方材料提供日期锚点，不能单独推知具体政策执行。这一变化若进入区域生活，必须通过道路、文书、军队、市场、家庭和地方中介才会发生；政权易手或法令发布不等于所有城镇、乡村与边地在同一时间获得同样经验。

本条所列来源为《元史》〈本纪〉及相关志；《元朝秘史》；《元典章》或《通制条格》相关条。它们能够支持特定的年序、制度语言、政治行动或局部遗存，却不能无条件化为社会总量。账本特别提醒“纪年框架可靠；该记录仅确证政权纪年连续，年内具体事项须由本年条另证”。因此，军报的兵数、条约的名分、官府的族群分类、判牍中的道德判断以及后出的叙述都须回到产生它们的场景，而不可直接代表全域动机或人口。

将事件与地方层次连接时，应同时考察资源怎样被征集、信息怎样传递、迁徙怎样改变户籍和劳动、宗教与亲属网络怎样提供或限制协作。现有证据往往只照亮少数地点；保留这种不均衡，才能避免把宋辽夏金元的多政权历史改写成单一中心的必然过程。

账本条目\texttt{YUAN-1207-REGNAL-YEAR}所记1207年的“蒙古以成吉思汗在位第2年作为诏令、赋役与外交文书的纪年框架”，属于蒙古的历史语境。条目所示前因是新岁颁历、年号和纪年是中枢与地方行政沟通的共同前提，其可见后果为为同年政令、战事和地方材料提供日期锚点，不能单独推知具体政策执行。这一变化若进入区域生活，必须通过道路、文书、军队、市场、家庭和地方中介才会发生；政权易手或法令发布不等于所有城镇、乡村与边地在同一时间获得同样经验。

本条所列来源为《元史》〈本纪〉及相关志；《元朝秘史》；《元典章》或《通制条格》相关条。它们能够支持特定的年序、制度语言、政治行动或局部遗存，却不能无条件化为社会总量。账本特别提醒“纪年框架可靠；该记录仅确证政权纪年连续，年内具体事项须由本年条另证”。因此，军报的兵数、条约的名分、官府的族群分类、判牍中的道德判断以及后出的叙述都须回到产生它们的场景，而不可直接代表全域动机或人口。

将事件与地方层次连接时，应同时考察资源怎样被征集、信息怎样传递、迁徙怎样改变户籍和劳动、宗教与亲属网络怎样提供或限制协作。现有证据往往只照亮少数地点；保留这种不均衡，才能避免把宋辽夏金元的多政权历史改写成单一中心的必然过程。

账本条目\texttt{YUAN-1208-REGNAL-YEAR}所记1208年的“蒙古以成吉思汗在位第3年作为诏令、赋役与外交文书的纪年框架”，属于蒙古的历史语境。条目所示前因是新岁颁历、年号和纪年是中枢与地方行政沟通的共同前提，其可见后果为为同年政令、战事和地方材料提供日期锚点，不能单独推知具体政策执行。这一变化若进入区域生活，必须通过道路、文书、军队、市场、家庭和地方中介才会发生；政权易手或法令发布不等于所有城镇、乡村与边地在同一时间获得同样经验。

本条所列来源为《元史》〈本纪〉及相关志；《元朝秘史》；《元典章》或《通制条格》相关条。它们能够支持特定的年序、制度语言、政治行动或局部遗存，却不能无条件化为社会总量。账本特别提醒“纪年框架可靠；该记录仅确证政权纪年连续，年内具体事项须由本年条另证”。因此，军报的兵数、条约的名分、官府的族群分类、判牍中的道德判断以及后出的叙述都须回到产生它们的场景，而不可直接代表全域动机或人口。

将事件与地方层次连接时，应同时考察资源怎样被征集、信息怎样传递、迁徙怎样改变户籍和劳动、宗教与亲属网络怎样提供或限制协作。现有证据往往只照亮少数地点；保留这种不均衡，才能避免把宋辽夏金元的多政权历史改写成单一中心的必然过程。

账本条目\texttt{YUAN-1209-EVENT-152}所记1209年的“蒙古军攻西夏并迫使其接受和议与贡赋”，属于蒙古与西夏的历史语境。条目所示前因是前期边防、财政和统治联盟的压力构成直接背景，具体条件须按本条材料分辨，其可见后果为改变后续资源配置与政治选择；实际执行范围和社会影响仍须分项核验。这一变化若进入区域生活，必须通过道路、文书、军队、市场、家庭和地方中介才会发生；政权易手或法令发布不等于所有城镇、乡村与边地在同一时间获得同样经验。

本条所列来源为《元史》〈本纪〉及相关志；《元朝秘史》；《元典章》或《通制条格》相关条。它们能够支持特定的年序、制度语言、政治行动或局部遗存，却不能无条件化为社会总量。账本特别提醒“编年主干可靠；细节、规模与动机须与对方材料、地方文书或考古材料互校”。因此，军报的兵数、条约的名分、官府的族群分类、判牍中的道德判断以及后出的叙述都须回到产生它们的场景，而不可直接代表全域动机或人口。

将事件与地方层次连接时，应同时考察资源怎样被征集、信息怎样传递、迁徙怎样改变户籍和劳动、宗教与亲属网络怎样提供或限制协作。现有证据往往只照亮少数地点；保留这种不均衡，才能避免把宋辽夏金元的多政权历史改写成单一中心的必然过程。

账本条目\texttt{YUAN-1209-REGNAL-YEAR}所记1209年的“蒙古以成吉思汗在位第4年作为诏令、赋役与外交文书的纪年框架”，属于蒙古的历史语境。条目所示前因是新岁颁历、年号和纪年是中枢与地方行政沟通的共同前提，其可见后果为为同年政令、战事和地方材料提供日期锚点，不能单独推知具体政策执行。这一变化若进入区域生活，必须通过道路、文书、军队、市场、家庭和地方中介才会发生；政权易手或法令发布不等于所有城镇、乡村与边地在同一时间获得同样经验。

本条所列来源为《元史》〈本纪〉及相关志；《元朝秘史》；《元典章》或《通制条格》相关条。它们能够支持特定的年序、制度语言、政治行动或局部遗存，却不能无条件化为社会总量。账本特别提醒“纪年框架可靠；该记录仅确证政权纪年连续，年内具体事项须由本年条另证”。因此，军报的兵数、条约的名分、官府的族群分类、判牍中的道德判断以及后出的叙述都须回到产生它们的场景，而不可直接代表全域动机或人口。

将事件与地方层次连接时，应同时考察资源怎样被征集、信息怎样传递、迁徙怎样改变户籍和劳动、宗教与亲属网络怎样提供或限制协作。现有证据往往只照亮少数地点；保留这种不均衡，才能避免把宋辽夏金元的多政权历史改写成单一中心的必然过程。

账本条目\texttt{YUAN-1210-REGNAL-YEAR}所记1210年的“蒙古以成吉思汗在位第5年作为诏令、赋役与外交文书的纪年框架”，属于蒙古的历史语境。条目所示前因是新岁颁历、年号和纪年是中枢与地方行政沟通的共同前提，其可见后果为为同年政令、战事和地方材料提供日期锚点，不能单独推知具体政策执行。这一变化若进入区域生活，必须通过道路、文书、军队、市场、家庭和地方中介才会发生；政权易手或法令发布不等于所有城镇、乡村与边地在同一时间获得同样经验。

本条所列来源为《元史》〈本纪〉及相关志；《元朝秘史》；《元典章》或《通制条格》相关条。它们能够支持特定的年序、制度语言、政治行动或局部遗存，却不能无条件化为社会总量。账本特别提醒“纪年框架可靠；该记录仅确证政权纪年连续，年内具体事项须由本年条另证”。因此，军报的兵数、条约的名分、官府的族群分类、判牍中的道德判断以及后出的叙述都须回到产生它们的场景，而不可直接代表全域动机或人口。

将事件与地方层次连接时，应同时考察资源怎样被征集、信息怎样传递、迁徙怎样改变户籍和劳动、宗教与亲属网络怎样提供或限制协作。现有证据往往只照亮少数地点；保留这种不均衡，才能避免把宋辽夏金元的多政权历史改写成单一中心的必然过程。

账本条目\texttt{YUAN-1211-EVENT-153}所记1211年的“蒙古大举进攻金朝北方边境”，属于蒙古与金的历史语境。条目所示前因是前期边防、财政和统治联盟的压力构成直接背景，具体条件须按本条材料分辨，其可见后果为改变后续资源配置与政治选择；实际执行范围和社会影响仍须分项核验。这一变化若进入区域生活，必须通过道路、文书、军队、市场、家庭和地方中介才会发生；政权易手或法令发布不等于所有城镇、乡村与边地在同一时间获得同样经验。

本条所列来源为《元史》〈本纪〉及相关志；《元朝秘史》；《元典章》或《通制条格》相关条。它们能够支持特定的年序、制度语言、政治行动或局部遗存，却不能无条件化为社会总量。账本特别提醒“编年主干可靠；细节、规模与动机须与对方材料、地方文书或考古材料互校”。因此，军报的兵数、条约的名分、官府的族群分类、判牍中的道德判断以及后出的叙述都须回到产生它们的场景，而不可直接代表全域动机或人口。

将事件与地方层次连接时，应同时考察资源怎样被征集、信息怎样传递、迁徙怎样改变户籍和劳动、宗教与亲属网络怎样提供或限制协作。现有证据往往只照亮少数地点；保留这种不均衡，才能避免把宋辽夏金元的多政权历史改写成单一中心的必然过程。

账本条目\texttt{YUAN-1211-REGNAL-YEAR}所记1211年的“蒙古以成吉思汗在位第6年作为诏令、赋役与外交文书的纪年框架”，属于蒙古的历史语境。条目所示前因是新岁颁历、年号和纪年是中枢与地方行政沟通的共同前提，其可见后果为为同年政令、战事和地方材料提供日期锚点，不能单独推知具体政策执行。这一变化若进入区域生活，必须通过道路、文书、军队、市场、家庭和地方中介才会发生；政权易手或法令发布不等于所有城镇、乡村与边地在同一时间获得同样经验。

本条所列来源为《元史》〈本纪〉及相关志；《元朝秘史》；《元典章》或《通制条格》相关条。它们能够支持特定的年序、制度语言、政治行动或局部遗存，却不能无条件化为社会总量。账本特别提醒“纪年框架可靠；该记录仅确证政权纪年连续，年内具体事项须由本年条另证”。因此，军报的兵数、条约的名分、官府的族群分类、判牍中的道德判断以及后出的叙述都须回到产生它们的场景，而不可直接代表全域动机或人口。

将事件与地方层次连接时，应同时考察资源怎样被征集、信息怎样传递、迁徙怎样改变户籍和劳动、宗教与亲属网络怎样提供或限制协作。现有证据往往只照亮少数地点；保留这种不均衡，才能避免把宋辽夏金元的多政权历史改写成单一中心的必然过程。

账本条目\texttt{YUAN-1212-REGNAL-YEAR}所记1212年的“蒙古以成吉思汗在位第7年作为诏令、赋役与外交文书的纪年框架”，属于蒙古的历史语境。条目所示前因是新岁颁历、年号和纪年是中枢与地方行政沟通的共同前提，其可见后果为为同年政令、战事和地方材料提供日期锚点，不能单独推知具体政策执行。这一变化若进入区域生活，必须通过道路、文书、军队、市场、家庭和地方中介才会发生；政权易手或法令发布不等于所有城镇、乡村与边地在同一时间获得同样经验。

本条所列来源为《元史》〈本纪〉及相关志；《元朝秘史》；《元典章》或《通制条格》相关条。它们能够支持特定的年序、制度语言、政治行动或局部遗存，却不能无条件化为社会总量。账本特别提醒“纪年框架可靠；该记录仅确证政权纪年连续，年内具体事项须由本年条另证”。因此，军报的兵数、条约的名分、官府的族群分类、判牍中的道德判断以及后出的叙述都须回到产生它们的场景，而不可直接代表全域动机或人口。

将事件与地方层次连接时，应同时考察资源怎样被征集、信息怎样传递、迁徙怎样改变户籍和劳动、宗教与亲属网络怎样提供或限制协作。现有证据往往只照亮少数地点；保留这种不均衡，才能避免把宋辽夏金元的多政权历史改写成单一中心的必然过程。

账本条目\texttt{YUAN-1213-REGNAL-YEAR}所记1213年的“蒙古以成吉思汗在位第8年作为诏令、赋役与外交文书的纪年框架”，属于蒙古的历史语境。条目所示前因是新岁颁历、年号和纪年是中枢与地方行政沟通的共同前提，其可见后果为为同年政令、战事和地方材料提供日期锚点，不能单独推知具体政策执行。这一变化若进入区域生活，必须通过道路、文书、军队、市场、家庭和地方中介才会发生；政权易手或法令发布不等于所有城镇、乡村与边地在同一时间获得同样经验。

本条所列来源为《元史》〈本纪〉及相关志；《元朝秘史》；《元典章》或《通制条格》相关条。它们能够支持特定的年序、制度语言、政治行动或局部遗存，却不能无条件化为社会总量。账本特别提醒“纪年框架可靠；该记录仅确证政权纪年连续，年内具体事项须由本年条另证”。因此，军报的兵数、条约的名分、官府的族群分类、判牍中的道德判断以及后出的叙述都须回到产生它们的场景，而不可直接代表全域动机或人口。

将事件与地方层次连接时，应同时考察资源怎样被征集、信息怎样传递、迁徙怎样改变户籍和劳动、宗教与亲属网络怎样提供或限制协作。现有证据往往只照亮少数地点；保留这种不均衡，才能避免把宋辽夏金元的多政权历史改写成单一中心的必然过程。

账本条目\texttt{YUAN-1214-REGNAL-YEAR}所记1214年的“蒙古以成吉思汗在位第9年作为诏令、赋役与外交文书的纪年框架”，属于蒙古的历史语境。条目所示前因是新岁颁历、年号和纪年是中枢与地方行政沟通的共同前提，其可见后果为为同年政令、战事和地方材料提供日期锚点，不能单独推知具体政策执行。这一变化若进入区域生活，必须通过道路、文书、军队、市场、家庭和地方中介才会发生；政权易手或法令发布不等于所有城镇、乡村与边地在同一时间获得同样经验。

本条所列来源为《元史》〈本纪〉及相关志；《元朝秘史》；《元典章》或《通制条格》相关条。它们能够支持特定的年序、制度语言、政治行动或局部遗存，却不能无条件化为社会总量。账本特别提醒“纪年框架可靠；该记录仅确证政权纪年连续，年内具体事项须由本年条另证”。因此，军报的兵数、条约的名分、官府的族群分类、判牍中的道德判断以及后出的叙述都须回到产生它们的场景，而不可直接代表全域动机或人口。

将事件与地方层次连接时，应同时考察资源怎样被征集、信息怎样传递、迁徙怎样改变户籍和劳动、宗教与亲属网络怎样提供或限制协作。现有证据往往只照亮少数地点；保留这种不均衡，才能避免把宋辽夏金元的多政权历史改写成单一中心的必然过程。

\subsection{外交、道路与人口移动}

账本条目\texttt{YUAN-1215-EVENT-154}所记1215年的“蒙古军攻占中都”，属于蒙古与金的历史语境。条目所示前因是前期边防、财政和统治联盟的压力构成直接背景，具体条件须按本条材料分辨，其可见后果为改变后续资源配置与政治选择；实际执行范围和社会影响仍须分项核验。这一变化若进入区域生活，必须通过道路、文书、军队、市场、家庭和地方中介才会发生；政权易手或法令发布不等于所有城镇、乡村与边地在同一时间获得同样经验。

本条所列来源为《元史》〈本纪〉及相关志；《元朝秘史》；《元典章》或《通制条格》相关条。它们能够支持特定的年序、制度语言、政治行动或局部遗存，却不能无条件化为社会总量。账本特别提醒“编年主干可靠；细节、规模与动机须与对方材料、地方文书或考古材料互校”。因此，军报的兵数、条约的名分、官府的族群分类、判牍中的道德判断以及后出的叙述都须回到产生它们的场景，而不可直接代表全域动机或人口。

将事件与地方层次连接时，应同时考察资源怎样被征集、信息怎样传递、迁徙怎样改变户籍和劳动、宗教与亲属网络怎样提供或限制协作。现有证据往往只照亮少数地点；保留这种不均衡，才能避免把宋辽夏金元的多政权历史改写成单一中心的必然过程。

账本条目\texttt{YUAN-1215-REGNAL-YEAR}所记1215年的“蒙古以成吉思汗在位第10年作为诏令、赋役与外交文书的纪年框架”，属于蒙古的历史语境。条目所示前因是新岁颁历、年号和纪年是中枢与地方行政沟通的共同前提，其可见后果为为同年政令、战事和地方材料提供日期锚点，不能单独推知具体政策执行。这一变化若进入区域生活，必须通过道路、文书、军队、市场、家庭和地方中介才会发生；政权易手或法令发布不等于所有城镇、乡村与边地在同一时间获得同样经验。

本条所列来源为《元史》〈本纪〉及相关志；《元朝秘史》；《元典章》或《通制条格》相关条。它们能够支持特定的年序、制度语言、政治行动或局部遗存，却不能无条件化为社会总量。账本特别提醒“纪年框架可靠；该记录仅确证政权纪年连续，年内具体事项须由本年条另证”。因此，军报的兵数、条约的名分、官府的族群分类、判牍中的道德判断以及后出的叙述都须回到产生它们的场景，而不可直接代表全域动机或人口。

将事件与地方层次连接时，应同时考察资源怎样被征集、信息怎样传递、迁徙怎样改变户籍和劳动、宗教与亲属网络怎样提供或限制协作。现有证据往往只照亮少数地点；保留这种不均衡，才能避免把宋辽夏金元的多政权历史改写成单一中心的必然过程。

账本条目\texttt{YUAN-1216-REGNAL-YEAR}所记1216年的“蒙古以成吉思汗在位第11年作为诏令、赋役与外交文书的纪年框架”，属于蒙古的历史语境。条目所示前因是新岁颁历、年号和纪年是中枢与地方行政沟通的共同前提，其可见后果为为同年政令、战事和地方材料提供日期锚点，不能单独推知具体政策执行。这一变化若进入区域生活，必须通过道路、文书、军队、市场、家庭和地方中介才会发生；政权易手或法令发布不等于所有城镇、乡村与边地在同一时间获得同样经验。

本条所列来源为《元史》〈本纪〉及相关志；《元朝秘史》；《元典章》或《通制条格》相关条。它们能够支持特定的年序、制度语言、政治行动或局部遗存，却不能无条件化为社会总量。账本特别提醒“纪年框架可靠；该记录仅确证政权纪年连续，年内具体事项须由本年条另证”。因此，军报的兵数、条约的名分、官府的族群分类、判牍中的道德判断以及后出的叙述都须回到产生它们的场景，而不可直接代表全域动机或人口。

将事件与地方层次连接时，应同时考察资源怎样被征集、信息怎样传递、迁徙怎样改变户籍和劳动、宗教与亲属网络怎样提供或限制协作。现有证据往往只照亮少数地点；保留这种不均衡，才能避免把宋辽夏金元的多政权历史改写成单一中心的必然过程。

账本条目\texttt{YUAN-1217-REGNAL-YEAR}所记1217年的“蒙古以成吉思汗在位第12年作为诏令、赋役与外交文书的纪年框架”，属于蒙古的历史语境。条目所示前因是新岁颁历、年号和纪年是中枢与地方行政沟通的共同前提，其可见后果为为同年政令、战事和地方材料提供日期锚点，不能单独推知具体政策执行。这一变化若进入区域生活，必须通过道路、文书、军队、市场、家庭和地方中介才会发生；政权易手或法令发布不等于所有城镇、乡村与边地在同一时间获得同样经验。

本条所列来源为《元史》〈本纪〉及相关志；《元朝秘史》；《元典章》或《通制条格》相关条。它们能够支持特定的年序、制度语言、政治行动或局部遗存，却不能无条件化为社会总量。账本特别提醒“纪年框架可靠；该记录仅确证政权纪年连续，年内具体事项须由本年条另证”。因此，军报的兵数、条约的名分、官府的族群分类、判牍中的道德判断以及后出的叙述都须回到产生它们的场景，而不可直接代表全域动机或人口。

将事件与地方层次连接时，应同时考察资源怎样被征集、信息怎样传递、迁徙怎样改变户籍和劳动、宗教与亲属网络怎样提供或限制协作。现有证据往往只照亮少数地点；保留这种不均衡，才能避免把宋辽夏金元的多政权历史改写成单一中心的必然过程。

账本条目\texttt{YUAN-1218-EVENT-155}所记1218年的“蒙古消灭西辽并进入中亚政治网络”，属于蒙古与西辽的历史语境。条目所示前因是前期边防、财政和统治联盟的压力构成直接背景，具体条件须按本条材料分辨，其可见后果为改变后续资源配置与政治选择；实际执行范围和社会影响仍须分项核验。这一变化若进入区域生活，必须通过道路、文书、军队、市场、家庭和地方中介才会发生；政权易手或法令发布不等于所有城镇、乡村与边地在同一时间获得同样经验。

本条所列来源为《元史》〈本纪〉及相关志；《元朝秘史》；《元典章》或《通制条格》相关条。它们能够支持特定的年序、制度语言、政治行动或局部遗存，却不能无条件化为社会总量。账本特别提醒“编年主干可靠；细节、规模与动机须与对方材料、地方文书或考古材料互校”。因此，军报的兵数、条约的名分、官府的族群分类、判牍中的道德判断以及后出的叙述都须回到产生它们的场景，而不可直接代表全域动机或人口。

将事件与地方层次连接时，应同时考察资源怎样被征集、信息怎样传递、迁徙怎样改变户籍和劳动、宗教与亲属网络怎样提供或限制协作。现有证据往往只照亮少数地点；保留这种不均衡，才能避免把宋辽夏金元的多政权历史改写成单一中心的必然过程。

账本条目\texttt{YUAN-1218-REGNAL-YEAR}所记1218年的“蒙古以成吉思汗在位第13年作为诏令、赋役与外交文书的纪年框架”，属于蒙古的历史语境。条目所示前因是新岁颁历、年号和纪年是中枢与地方行政沟通的共同前提，其可见后果为为同年政令、战事和地方材料提供日期锚点，不能单独推知具体政策执行。这一变化若进入区域生活，必须通过道路、文书、军队、市场、家庭和地方中介才会发生；政权易手或法令发布不等于所有城镇、乡村与边地在同一时间获得同样经验。

本条所列来源为《元史》〈本纪〉及相关志；《元朝秘史》；《元典章》或《通制条格》相关条。它们能够支持特定的年序、制度语言、政治行动或局部遗存，却不能无条件化为社会总量。账本特别提醒“纪年框架可靠；该记录仅确证政权纪年连续，年内具体事项须由本年条另证”。因此，军报的兵数、条约的名分、官府的族群分类、判牍中的道德判断以及后出的叙述都须回到产生它们的场景，而不可直接代表全域动机或人口。

将事件与地方层次连接时，应同时考察资源怎样被征集、信息怎样传递、迁徙怎样改变户籍和劳动、宗教与亲属网络怎样提供或限制协作。现有证据往往只照亮少数地点；保留这种不均衡，才能避免把宋辽夏金元的多政权历史改写成单一中心的必然过程。

账本条目\texttt{YUAN-1219-EVENT-156}所记1219年的“蒙古发动对花剌子模的大规模战争”，属于蒙古与花剌子模的历史语境。条目所示前因是前期边防、财政和统治联盟的压力构成直接背景，具体条件须按本条材料分辨，其可见后果为改变后续资源配置与政治选择；实际执行范围和社会影响仍须分项核验。这一变化若进入区域生活，必须通过道路、文书、军队、市场、家庭和地方中介才会发生；政权易手或法令发布不等于所有城镇、乡村与边地在同一时间获得同样经验。

本条所列来源为《元史》〈本纪〉及相关志；《元朝秘史》；《元典章》或《通制条格》相关条。它们能够支持特定的年序、制度语言、政治行动或局部遗存，却不能无条件化为社会总量。账本特别提醒“编年主干可靠；细节、规模与动机须与对方材料、地方文书或考古材料互校”。因此，军报的兵数、条约的名分、官府的族群分类、判牍中的道德判断以及后出的叙述都须回到产生它们的场景，而不可直接代表全域动机或人口。

将事件与地方层次连接时，应同时考察资源怎样被征集、信息怎样传递、迁徙怎样改变户籍和劳动、宗教与亲属网络怎样提供或限制协作。现有证据往往只照亮少数地点；保留这种不均衡，才能避免把宋辽夏金元的多政权历史改写成单一中心的必然过程。

账本条目\texttt{YUAN-1219-REGNAL-YEAR}所记1219年的“蒙古以成吉思汗在位第14年作为诏令、赋役与外交文书的纪年框架”，属于蒙古的历史语境。条目所示前因是新岁颁历、年号和纪年是中枢与地方行政沟通的共同前提，其可见后果为为同年政令、战事和地方材料提供日期锚点，不能单独推知具体政策执行。这一变化若进入区域生活，必须通过道路、文书、军队、市场、家庭和地方中介才会发生；政权易手或法令发布不等于所有城镇、乡村与边地在同一时间获得同样经验。

本条所列来源为《元史》〈本纪〉及相关志；《元朝秘史》；《元典章》或《通制条格》相关条。它们能够支持特定的年序、制度语言、政治行动或局部遗存，却不能无条件化为社会总量。账本特别提醒“纪年框架可靠；该记录仅确证政权纪年连续，年内具体事项须由本年条另证”。因此，军报的兵数、条约的名分、官府的族群分类、判牍中的道德判断以及后出的叙述都须回到产生它们的场景，而不可直接代表全域动机或人口。

将事件与地方层次连接时，应同时考察资源怎样被征集、信息怎样传递、迁徙怎样改变户籍和劳动、宗教与亲属网络怎样提供或限制协作。现有证据往往只照亮少数地点；保留这种不均衡，才能避免把宋辽夏金元的多政权历史改写成单一中心的必然过程。

账本条目\texttt{YUAN-1220-REGNAL-YEAR}所记1220年的“蒙古以成吉思汗在位第15年作为诏令、赋役与外交文书的纪年框架”，属于蒙古的历史语境。条目所示前因是新岁颁历、年号和纪年是中枢与地方行政沟通的共同前提，其可见后果为为同年政令、战事和地方材料提供日期锚点，不能单独推知具体政策执行。这一变化若进入区域生活，必须通过道路、文书、军队、市场、家庭和地方中介才会发生；政权易手或法令发布不等于所有城镇、乡村与边地在同一时间获得同样经验。

本条所列来源为《元史》〈本纪〉及相关志；《元朝秘史》；《元典章》或《通制条格》相关条。它们能够支持特定的年序、制度语言、政治行动或局部遗存，却不能无条件化为社会总量。账本特别提醒“纪年框架可靠；该记录仅确证政权纪年连续，年内具体事项须由本年条另证”。因此，军报的兵数、条约的名分、官府的族群分类、判牍中的道德判断以及后出的叙述都须回到产生它们的场景，而不可直接代表全域动机或人口。

将事件与地方层次连接时，应同时考察资源怎样被征集、信息怎样传递、迁徙怎样改变户籍和劳动、宗教与亲属网络怎样提供或限制协作。现有证据往往只照亮少数地点；保留这种不均衡，才能避免把宋辽夏金元的多政权历史改写成单一中心的必然过程。

账本条目\texttt{YUAN-1221-REGNAL-YEAR}所记1221年的“蒙古以成吉思汗在位第16年作为诏令、赋役与外交文书的纪年框架”，属于蒙古的历史语境。条目所示前因是新岁颁历、年号和纪年是中枢与地方行政沟通的共同前提，其可见后果为为同年政令、战事和地方材料提供日期锚点，不能单独推知具体政策执行。这一变化若进入区域生活，必须通过道路、文书、军队、市场、家庭和地方中介才会发生；政权易手或法令发布不等于所有城镇、乡村与边地在同一时间获得同样经验。

本条所列来源为《元史》〈本纪〉及相关志；《元朝秘史》；《元典章》或《通制条格》相关条。它们能够支持特定的年序、制度语言、政治行动或局部遗存，却不能无条件化为社会总量。账本特别提醒“纪年框架可靠；该记录仅确证政权纪年连续，年内具体事项须由本年条另证”。因此，军报的兵数、条约的名分、官府的族群分类、判牍中的道德判断以及后出的叙述都须回到产生它们的场景，而不可直接代表全域动机或人口。

将事件与地方层次连接时，应同时考察资源怎样被征集、信息怎样传递、迁徙怎样改变户籍和劳动、宗教与亲属网络怎样提供或限制协作。现有证据往往只照亮少数地点；保留这种不均衡，才能避免把宋辽夏金元的多政权历史改写成单一中心的必然过程。

账本条目\texttt{YUAN-1222-REGNAL-YEAR}所记1222年的“蒙古以成吉思汗在位第17年作为诏令、赋役与外交文书的纪年框架”，属于蒙古的历史语境。条目所示前因是新岁颁历、年号和纪年是中枢与地方行政沟通的共同前提，其可见后果为为同年政令、战事和地方材料提供日期锚点，不能单独推知具体政策执行。这一变化若进入区域生活，必须通过道路、文书、军队、市场、家庭和地方中介才会发生；政权易手或法令发布不等于所有城镇、乡村与边地在同一时间获得同样经验。

本条所列来源为《元史》〈本纪〉及相关志；《元朝秘史》；《元典章》或《通制条格》相关条。它们能够支持特定的年序、制度语言、政治行动或局部遗存，却不能无条件化为社会总量。账本特别提醒“纪年框架可靠；该记录仅确证政权纪年连续，年内具体事项须由本年条另证”。因此，军报的兵数、条约的名分、官府的族群分类、判牍中的道德判断以及后出的叙述都须回到产生它们的场景，而不可直接代表全域动机或人口。

将事件与地方层次连接时，应同时考察资源怎样被征集、信息怎样传递、迁徙怎样改变户籍和劳动、宗教与亲属网络怎样提供或限制协作。现有证据往往只照亮少数地点；保留这种不均衡，才能避免把宋辽夏金元的多政权历史改写成单一中心的必然过程。

账本条目\texttt{YUAN-1223-REGNAL-YEAR}所记1223年的“蒙古以成吉思汗在位第18年作为诏令、赋役与外交文书的纪年框架”，属于蒙古的历史语境。条目所示前因是新岁颁历、年号和纪年是中枢与地方行政沟通的共同前提，其可见后果为为同年政令、战事和地方材料提供日期锚点，不能单独推知具体政策执行。这一变化若进入区域生活，必须通过道路、文书、军队、市场、家庭和地方中介才会发生；政权易手或法令发布不等于所有城镇、乡村与边地在同一时间获得同样经验。

本条所列来源为《元史》〈本纪〉及相关志；《元朝秘史》；《元典章》或《通制条格》相关条。它们能够支持特定的年序、制度语言、政治行动或局部遗存，却不能无条件化为社会总量。账本特别提醒“纪年框架可靠；该记录仅确证政权纪年连续，年内具体事项须由本年条另证”。因此，军报的兵数、条约的名分、官府的族群分类、判牍中的道德判断以及后出的叙述都须回到产生它们的场景，而不可直接代表全域动机或人口。

将事件与地方层次连接时，应同时考察资源怎样被征集、信息怎样传递、迁徙怎样改变户籍和劳动、宗教与亲属网络怎样提供或限制协作。现有证据往往只照亮少数地点；保留这种不均衡，才能避免把宋辽夏金元的多政权历史改写成单一中心的必然过程。

\subsection{环境风险与地方经济}

账本条目\texttt{YUAN-1224-REGNAL-YEAR}所记1224年的“蒙古以成吉思汗在位第19年作为诏令、赋役与外交文书的纪年框架”，属于蒙古的历史语境。条目所示前因是新岁颁历、年号和纪年是中枢与地方行政沟通的共同前提，其可见后果为为同年政令、战事和地方材料提供日期锚点，不能单独推知具体政策执行。这一变化若进入区域生活，必须通过道路、文书、军队、市场、家庭和地方中介才会发生；政权易手或法令发布不等于所有城镇、乡村与边地在同一时间获得同样经验。

本条所列来源为《元史》〈本纪〉及相关志；《元朝秘史》；《元典章》或《通制条格》相关条。它们能够支持特定的年序、制度语言、政治行动或局部遗存，却不能无条件化为社会总量。账本特别提醒“纪年框架可靠；该记录仅确证政权纪年连续，年内具体事项须由本年条另证”。因此，军报的兵数、条约的名分、官府的族群分类、判牍中的道德判断以及后出的叙述都须回到产生它们的场景，而不可直接代表全域动机或人口。

将事件与地方层次连接时，应同时考察资源怎样被征集、信息怎样传递、迁徙怎样改变户籍和劳动、宗教与亲属网络怎样提供或限制协作。现有证据往往只照亮少数地点；保留这种不均衡，才能避免把宋辽夏金元的多政权历史改写成单一中心的必然过程。

账本条目\texttt{YUAN-1225-REGNAL-YEAR}所记1225年的“蒙古以成吉思汗在位第20年作为诏令、赋役与外交文书的纪年框架”，属于蒙古的历史语境。条目所示前因是新岁颁历、年号和纪年是中枢与地方行政沟通的共同前提，其可见后果为为同年政令、战事和地方材料提供日期锚点，不能单独推知具体政策执行。这一变化若进入区域生活，必须通过道路、文书、军队、市场、家庭和地方中介才会发生；政权易手或法令发布不等于所有城镇、乡村与边地在同一时间获得同样经验。

本条所列来源为《元史》〈本纪〉及相关志；《元朝秘史》；《元典章》或《通制条格》相关条。它们能够支持特定的年序、制度语言、政治行动或局部遗存，却不能无条件化为社会总量。账本特别提醒“纪年框架可靠；该记录仅确证政权纪年连续，年内具体事项须由本年条另证”。因此，军报的兵数、条约的名分、官府的族群分类、判牍中的道德判断以及后出的叙述都须回到产生它们的场景，而不可直接代表全域动机或人口。

将事件与地方层次连接时，应同时考察资源怎样被征集、信息怎样传递、迁徙怎样改变户籍和劳动、宗教与亲属网络怎样提供或限制协作。现有证据往往只照亮少数地点；保留这种不均衡，才能避免把宋辽夏金元的多政权历史改写成单一中心的必然过程。

账本条目\texttt{YUAN-1226-REGNAL-YEAR}所记1226年的“蒙古以成吉思汗在位第21年作为诏令、赋役与外交文书的纪年框架”，属于蒙古的历史语境。条目所示前因是新岁颁历、年号和纪年是中枢与地方行政沟通的共同前提，其可见后果为为同年政令、战事和地方材料提供日期锚点，不能单独推知具体政策执行。这一变化若进入区域生活，必须通过道路、文书、军队、市场、家庭和地方中介才会发生；政权易手或法令发布不等于所有城镇、乡村与边地在同一时间获得同样经验。

本条所列来源为《元史》〈本纪〉及相关志；《元朝秘史》；《元典章》或《通制条格》相关条。它们能够支持特定的年序、制度语言、政治行动或局部遗存，却不能无条件化为社会总量。账本特别提醒“纪年框架可靠；该记录仅确证政权纪年连续，年内具体事项须由本年条另证”。因此，军报的兵数、条约的名分、官府的族群分类、判牍中的道德判断以及后出的叙述都须回到产生它们的场景，而不可直接代表全域动机或人口。

将事件与地方层次连接时，应同时考察资源怎样被征集、信息怎样传递、迁徙怎样改变户籍和劳动、宗教与亲属网络怎样提供或限制协作。现有证据往往只照亮少数地点；保留这种不均衡，才能避免把宋辽夏金元的多政权历史改写成单一中心的必然过程。

账本条目\texttt{YUAN-1227-EVENT-157}所记1227年的“西夏覆亡，同年成吉思汗去世”，属于蒙古与西夏的历史语境。条目所示前因是前期边防、财政和统治联盟的压力构成直接背景，具体条件须按本条材料分辨，其可见后果为改变后续资源配置与政治选择；实际执行范围和社会影响仍须分项核验。这一变化若进入区域生活，必须通过道路、文书、军队、市场、家庭和地方中介才会发生；政权易手或法令发布不等于所有城镇、乡村与边地在同一时间获得同样经验。

本条所列来源为《元史》〈本纪〉及相关志；《元朝秘史》；《元典章》或《通制条格》相关条。它们能够支持特定的年序、制度语言、政治行动或局部遗存，却不能无条件化为社会总量。账本特别提醒“编年主干可靠；细节、规模与动机须与对方材料、地方文书或考古材料互校”。因此，军报的兵数、条约的名分、官府的族群分类、判牍中的道德判断以及后出的叙述都须回到产生它们的场景，而不可直接代表全域动机或人口。

将事件与地方层次连接时，应同时考察资源怎样被征集、信息怎样传递、迁徙怎样改变户籍和劳动、宗教与亲属网络怎样提供或限制协作。现有证据往往只照亮少数地点；保留这种不均衡，才能避免把宋辽夏金元的多政权历史改写成单一中心的必然过程。

账本条目\texttt{YUAN-1227-REGNAL-YEAR}所记1227年的“蒙古以成吉思汗在位第22年作为诏令、赋役与外交文书的纪年框架”，属于蒙古的历史语境。条目所示前因是新岁颁历、年号和纪年是中枢与地方行政沟通的共同前提，其可见后果为为同年政令、战事和地方材料提供日期锚点，不能单独推知具体政策执行。这一变化若进入区域生活，必须通过道路、文书、军队、市场、家庭和地方中介才会发生；政权易手或法令发布不等于所有城镇、乡村与边地在同一时间获得同样经验。

本条所列来源为《元史》〈本纪〉及相关志；《元朝秘史》；《元典章》或《通制条格》相关条。它们能够支持特定的年序、制度语言、政治行动或局部遗存，却不能无条件化为社会总量。账本特别提醒“纪年框架可靠；该记录仅确证政权纪年连续，年内具体事项须由本年条另证”。因此，军报的兵数、条约的名分、官府的族群分类、判牍中的道德判断以及后出的叙述都须回到产生它们的场景，而不可直接代表全域动机或人口。

将事件与地方层次连接时，应同时考察资源怎样被征集、信息怎样传递、迁徙怎样改变户籍和劳动、宗教与亲属网络怎样提供或限制协作。现有证据往往只照亮少数地点；保留这种不均衡，才能避免把宋辽夏金元的多政权历史改写成单一中心的必然过程。

账本条目\texttt{YUAN-1228-REGNAL-YEAR}所记1228年的“蒙古以窝阔台在位第1年作为诏令、赋役与外交文书的纪年框架”，属于蒙古的历史语境。条目所示前因是新岁颁历、年号和纪年是中枢与地方行政沟通的共同前提，其可见后果为为同年政令、战事和地方材料提供日期锚点，不能单独推知具体政策执行。这一变化若进入区域生活，必须通过道路、文书、军队、市场、家庭和地方中介才会发生；政权易手或法令发布不等于所有城镇、乡村与边地在同一时间获得同样经验。

本条所列来源为《元史》〈本纪〉及相关志；《元朝秘史》；《元典章》或《通制条格》相关条。它们能够支持特定的年序、制度语言、政治行动或局部遗存，却不能无条件化为社会总量。账本特别提醒“纪年框架可靠；该记录仅确证政权纪年连续，年内具体事项须由本年条另证”。因此，军报的兵数、条约的名分、官府的族群分类、判牍中的道德判断以及后出的叙述都须回到产生它们的场景，而不可直接代表全域动机或人口。

将事件与地方层次连接时，应同时考察资源怎样被征集、信息怎样传递、迁徙怎样改变户籍和劳动、宗教与亲属网络怎样提供或限制协作。现有证据往往只照亮少数地点；保留这种不均衡，才能避免把宋辽夏金元的多政权历史改写成单一中心的必然过程。

账本条目\texttt{YUAN-1229-EVENT-158}所记1229年的“窝阔台即大汗位”，属于蒙古的历史语境。条目所示前因是前期边防、财政和统治联盟的压力构成直接背景，具体条件须按本条材料分辨，其可见后果为改变后续资源配置与政治选择；实际执行范围和社会影响仍须分项核验。这一变化若进入区域生活，必须通过道路、文书、军队、市场、家庭和地方中介才会发生；政权易手或法令发布不等于所有城镇、乡村与边地在同一时间获得同样经验。

本条所列来源为《元史》〈本纪〉及相关志；《元朝秘史》；《元典章》或《通制条格》相关条。它们能够支持特定的年序、制度语言、政治行动或局部遗存，却不能无条件化为社会总量。账本特别提醒“编年主干可靠；细节、规模与动机须与对方材料、地方文书或考古材料互校”。因此，军报的兵数、条约的名分、官府的族群分类、判牍中的道德判断以及后出的叙述都须回到产生它们的场景，而不可直接代表全域动机或人口。

将事件与地方层次连接时，应同时考察资源怎样被征集、信息怎样传递、迁徙怎样改变户籍和劳动、宗教与亲属网络怎样提供或限制协作。现有证据往往只照亮少数地点；保留这种不均衡，才能避免把宋辽夏金元的多政权历史改写成单一中心的必然过程。

账本条目\texttt{YUAN-1229-REGNAL-YEAR}所记1229年的“蒙古以窝阔台在位第2年作为诏令、赋役与外交文书的纪年框架”，属于蒙古的历史语境。条目所示前因是新岁颁历、年号和纪年是中枢与地方行政沟通的共同前提，其可见后果为为同年政令、战事和地方材料提供日期锚点，不能单独推知具体政策执行。这一变化若进入区域生活，必须通过道路、文书、军队、市场、家庭和地方中介才会发生；政权易手或法令发布不等于所有城镇、乡村与边地在同一时间获得同样经验。

本条所列来源为《元史》〈本纪〉及相关志；《元朝秘史》；《元典章》或《通制条格》相关条。它们能够支持特定的年序、制度语言、政治行动或局部遗存，却不能无条件化为社会总量。账本特别提醒“纪年框架可靠；该记录仅确证政权纪年连续，年内具体事项须由本年条另证”。因此，军报的兵数、条约的名分、官府的族群分类、判牍中的道德判断以及后出的叙述都须回到产生它们的场景，而不可直接代表全域动机或人口。

将事件与地方层次连接时，应同时考察资源怎样被征集、信息怎样传递、迁徙怎样改变户籍和劳动、宗教与亲属网络怎样提供或限制协作。现有证据往往只照亮少数地点；保留这种不均衡，才能避免把宋辽夏金元的多政权历史改写成单一中心的必然过程。

账本条目\texttt{YUAN-1230-REGNAL-YEAR}所记1230年的“蒙古以窝阔台在位第3年作为诏令、赋役与外交文书的纪年框架”，属于蒙古的历史语境。条目所示前因是新岁颁历、年号和纪年是中枢与地方行政沟通的共同前提，其可见后果为为同年政令、战事和地方材料提供日期锚点，不能单独推知具体政策执行。这一变化若进入区域生活，必须通过道路、文书、军队、市场、家庭和地方中介才会发生；政权易手或法令发布不等于所有城镇、乡村与边地在同一时间获得同样经验。

本条所列来源为《元史》〈本纪〉及相关志；《元朝秘史》；《元典章》或《通制条格》相关条。它们能够支持特定的年序、制度语言、政治行动或局部遗存，却不能无条件化为社会总量。账本特别提醒“纪年框架可靠；该记录仅确证政权纪年连续，年内具体事项须由本年条另证”。因此，军报的兵数、条约的名分、官府的族群分类、判牍中的道德判断以及后出的叙述都须回到产生它们的场景，而不可直接代表全域动机或人口。

将事件与地方层次连接时，应同时考察资源怎样被征集、信息怎样传递、迁徙怎样改变户籍和劳动、宗教与亲属网络怎样提供或限制协作。现有证据往往只照亮少数地点；保留这种不均衡，才能避免把宋辽夏金元的多政权历史改写成单一中心的必然过程。

账本条目\texttt{YUAN-1231-EVENT-159}所记1231年的“三峰山等战事削弱金军主力”，属于蒙古与金的历史语境。条目所示前因是前期边防、财政和统治联盟的压力构成直接背景，具体条件须按本条材料分辨，其可见后果为改变后续资源配置与政治选择；实际执行范围和社会影响仍须分项核验。这一变化若进入区域生活，必须通过道路、文书、军队、市场、家庭和地方中介才会发生；政权易手或法令发布不等于所有城镇、乡村与边地在同一时间获得同样经验。

本条所列来源为《元史》〈本纪〉及相关志；《元朝秘史》；《元典章》或《通制条格》相关条。它们能够支持特定的年序、制度语言、政治行动或局部遗存，却不能无条件化为社会总量。账本特别提醒“编年主干可靠；细节、规模与动机须与对方材料、地方文书或考古材料互校”。因此，军报的兵数、条约的名分、官府的族群分类、判牍中的道德判断以及后出的叙述都须回到产生它们的场景，而不可直接代表全域动机或人口。

将事件与地方层次连接时，应同时考察资源怎样被征集、信息怎样传递、迁徙怎样改变户籍和劳动、宗教与亲属网络怎样提供或限制协作。现有证据往往只照亮少数地点；保留这种不均衡，才能避免把宋辽夏金元的多政权历史改写成单一中心的必然过程。

账本条目\texttt{YUAN-1231-REGNAL-YEAR}所记1231年的“蒙古以窝阔台在位第4年作为诏令、赋役与外交文书的纪年框架”，属于蒙古的历史语境。条目所示前因是新岁颁历、年号和纪年是中枢与地方行政沟通的共同前提，其可见后果为为同年政令、战事和地方材料提供日期锚点，不能单独推知具体政策执行。这一变化若进入区域生活，必须通过道路、文书、军队、市场、家庭和地方中介才会发生；政权易手或法令发布不等于所有城镇、乡村与边地在同一时间获得同样经验。

本条所列来源为《元史》〈本纪〉及相关志；《元朝秘史》；《元典章》或《通制条格》相关条。它们能够支持特定的年序、制度语言、政治行动或局部遗存，却不能无条件化为社会总量。账本特别提醒“纪年框架可靠；该记录仅确证政权纪年连续，年内具体事项须由本年条另证”。因此，军报的兵数、条约的名分、官府的族群分类、判牍中的道德判断以及后出的叙述都须回到产生它们的场景，而不可直接代表全域动机或人口。

将事件与地方层次连接时，应同时考察资源怎样被征集、信息怎样传递、迁徙怎样改变户籍和劳动、宗教与亲属网络怎样提供或限制协作。现有证据往往只照亮少数地点；保留这种不均衡，才能避免把宋辽夏金元的多政权历史改写成单一中心的必然过程。

账本条目\texttt{YUAN-1232-REGNAL-YEAR}所记1232年的“蒙古以窝阔台在位第5年作为诏令、赋役与外交文书的纪年框架”，属于蒙古的历史语境。条目所示前因是新岁颁历、年号和纪年是中枢与地方行政沟通的共同前提，其可见后果为为同年政令、战事和地方材料提供日期锚点，不能单独推知具体政策执行。这一变化若进入区域生活，必须通过道路、文书、军队、市场、家庭和地方中介才会发生；政权易手或法令发布不等于所有城镇、乡村与边地在同一时间获得同样经验。

本条所列来源为《元史》〈本纪〉及相关志；《元朝秘史》；《元典章》或《通制条格》相关条。它们能够支持特定的年序、制度语言、政治行动或局部遗存，却不能无条件化为社会总量。账本特别提醒“纪年框架可靠；该记录仅确证政权纪年连续，年内具体事项须由本年条另证”。因此，军报的兵数、条约的名分、官府的族群分类、判牍中的道德判断以及后出的叙述都须回到产生它们的场景，而不可直接代表全域动机或人口。

将事件与地方层次连接时，应同时考察资源怎样被征集、信息怎样传递、迁徙怎样改变户籍和劳动、宗教与亲属网络怎样提供或限制协作。现有证据往往只照亮少数地点；保留这种不均衡，才能避免把宋辽夏金元的多政权历史改写成单一中心的必然过程。

\subsection{宗教、法律与材料边界}

账本条目\texttt{YUAN-1233-REGNAL-YEAR}所记1233年的“蒙古以窝阔台在位第6年作为诏令、赋役与外交文书的纪年框架”，属于蒙古的历史语境。条目所示前因是新岁颁历、年号和纪年是中枢与地方行政沟通的共同前提，其可见后果为为同年政令、战事和地方材料提供日期锚点，不能单独推知具体政策执行。这一变化若进入区域生活，必须通过道路、文书、军队、市场、家庭和地方中介才会发生；政权易手或法令发布不等于所有城镇、乡村与边地在同一时间获得同样经验。

本条所列来源为《元史》〈本纪〉及相关志；《元朝秘史》；《元典章》或《通制条格》相关条。它们能够支持特定的年序、制度语言、政治行动或局部遗存，却不能无条件化为社会总量。账本特别提醒“纪年框架可靠；该记录仅确证政权纪年连续，年内具体事项须由本年条另证”。因此，军报的兵数、条约的名分、官府的族群分类、判牍中的道德判断以及后出的叙述都须回到产生它们的场景，而不可直接代表全域动机或人口。

将事件与地方层次连接时，应同时考察资源怎样被征集、信息怎样传递、迁徙怎样改变户籍和劳动、宗教与亲属网络怎样提供或限制协作。现有证据往往只照亮少数地点；保留这种不均衡，才能避免把宋辽夏金元的多政权历史改写成单一中心的必然过程。

账本条目\texttt{YUAN-1234-EVENT-160}所记1234年的“蒙古与南宋攻灭金朝”，属于蒙古南宋金的历史语境。条目所示前因是前期边防、财政和统治联盟的压力构成直接背景，具体条件须按本条材料分辨，其可见后果为改变后续资源配置与政治选择；实际执行范围和社会影响仍须分项核验。这一变化若进入区域生活，必须通过道路、文书、军队、市场、家庭和地方中介才会发生；政权易手或法令发布不等于所有城镇、乡村与边地在同一时间获得同样经验。

本条所列来源为《元史》〈本纪〉及相关志；《元朝秘史》；《元典章》或《通制条格》相关条。它们能够支持特定的年序、制度语言、政治行动或局部遗存，却不能无条件化为社会总量。账本特别提醒“编年主干可靠；细节、规模与动机须与对方材料、地方文书或考古材料互校”。因此，军报的兵数、条约的名分、官府的族群分类、判牍中的道德判断以及后出的叙述都须回到产生它们的场景，而不可直接代表全域动机或人口。

将事件与地方层次连接时，应同时考察资源怎样被征集、信息怎样传递、迁徙怎样改变户籍和劳动、宗教与亲属网络怎样提供或限制协作。现有证据往往只照亮少数地点；保留这种不均衡，才能避免把宋辽夏金元的多政权历史改写成单一中心的必然过程。

账本条目\texttt{YUAN-1234-REGNAL-YEAR}所记1234年的“蒙古以窝阔台在位第7年作为诏令、赋役与外交文书的纪年框架”，属于蒙古的历史语境。条目所示前因是新岁颁历、年号和纪年是中枢与地方行政沟通的共同前提，其可见后果为为同年政令、战事和地方材料提供日期锚点，不能单独推知具体政策执行。这一变化若进入区域生活，必须通过道路、文书、军队、市场、家庭和地方中介才会发生；政权易手或法令发布不等于所有城镇、乡村与边地在同一时间获得同样经验。

本条所列来源为《元史》〈本纪〉及相关志；《元朝秘史》；《元典章》或《通制条格》相关条。它们能够支持特定的年序、制度语言、政治行动或局部遗存，却不能无条件化为社会总量。账本特别提醒“纪年框架可靠；该记录仅确证政权纪年连续，年内具体事项须由本年条另证”。因此，军报的兵数、条约的名分、官府的族群分类、判牍中的道德判断以及后出的叙述都须回到产生它们的场景，而不可直接代表全域动机或人口。

将事件与地方层次连接时，应同时考察资源怎样被征集、信息怎样传递、迁徙怎样改变户籍和劳动、宗教与亲属网络怎样提供或限制协作。现有证据往往只照亮少数地点；保留这种不均衡，才能避免把宋辽夏金元的多政权历史改写成单一中心的必然过程。

账本条目\texttt{YUAN-1235-EVENT-161}所记1235年的“蒙古贵族会议决定多方向扩张”，属于蒙古的历史语境。条目所示前因是前期边防、财政和统治联盟的压力构成直接背景，具体条件须按本条材料分辨，其可见后果为改变后续资源配置与政治选择；实际执行范围和社会影响仍须分项核验。这一变化若进入区域生活，必须通过道路、文书、军队、市场、家庭和地方中介才会发生；政权易手或法令发布不等于所有城镇、乡村与边地在同一时间获得同样经验。

本条所列来源为《元史》〈本纪〉及相关志；《元朝秘史》；《元典章》或《通制条格》相关条。它们能够支持特定的年序、制度语言、政治行动或局部遗存，却不能无条件化为社会总量。账本特别提醒“编年主干可靠；细节、规模与动机须与对方材料、地方文书或考古材料互校”。因此，军报的兵数、条约的名分、官府的族群分类、判牍中的道德判断以及后出的叙述都须回到产生它们的场景，而不可直接代表全域动机或人口。

将事件与地方层次连接时，应同时考察资源怎样被征集、信息怎样传递、迁徙怎样改变户籍和劳动、宗教与亲属网络怎样提供或限制协作。现有证据往往只照亮少数地点；保留这种不均衡，才能避免把宋辽夏金元的多政权历史改写成单一中心的必然过程。

账本条目\texttt{YUAN-1235-REGNAL-YEAR}所记1235年的“蒙古以窝阔台在位第8年作为诏令、赋役与外交文书的纪年框架”，属于蒙古的历史语境。条目所示前因是新岁颁历、年号和纪年是中枢与地方行政沟通的共同前提，其可见后果为为同年政令、战事和地方材料提供日期锚点，不能单独推知具体政策执行。这一变化若进入区域生活，必须通过道路、文书、军队、市场、家庭和地方中介才会发生；政权易手或法令发布不等于所有城镇、乡村与边地在同一时间获得同样经验。

本条所列来源为《元史》〈本纪〉及相关志；《元朝秘史》；《元典章》或《通制条格》相关条。它们能够支持特定的年序、制度语言、政治行动或局部遗存，却不能无条件化为社会总量。账本特别提醒“纪年框架可靠；该记录仅确证政权纪年连续，年内具体事项须由本年条另证”。因此，军报的兵数、条约的名分、官府的族群分类、判牍中的道德判断以及后出的叙述都须回到产生它们的场景，而不可直接代表全域动机或人口。

将事件与地方层次连接时，应同时考察资源怎样被征集、信息怎样传递、迁徙怎样改变户籍和劳动、宗教与亲属网络怎样提供或限制协作。现有证据往往只照亮少数地点；保留这种不均衡，才能避免把宋辽夏金元的多政权历史改写成单一中心的必然过程。

账本条目\texttt{YUAN-1236-REGNAL-YEAR}所记1236年的“蒙古以窝阔台在位第9年作为诏令、赋役与外交文书的纪年框架”，属于蒙古的历史语境。条目所示前因是新岁颁历、年号和纪年是中枢与地方行政沟通的共同前提，其可见后果为为同年政令、战事和地方材料提供日期锚点，不能单独推知具体政策执行。这一变化若进入区域生活，必须通过道路、文书、军队、市场、家庭和地方中介才会发生；政权易手或法令发布不等于所有城镇、乡村与边地在同一时间获得同样经验。

本条所列来源为《元史》〈本纪〉及相关志；《元朝秘史》；《元典章》或《通制条格》相关条。它们能够支持特定的年序、制度语言、政治行动或局部遗存，却不能无条件化为社会总量。账本特别提醒“纪年框架可靠；该记录仅确证政权纪年连续，年内具体事项须由本年条另证”。因此，军报的兵数、条约的名分、官府的族群分类、判牍中的道德判断以及后出的叙述都须回到产生它们的场景，而不可直接代表全域动机或人口。

将事件与地方层次连接时，应同时考察资源怎样被征集、信息怎样传递、迁徙怎样改变户籍和劳动、宗教与亲属网络怎样提供或限制协作。现有证据往往只照亮少数地点；保留这种不均衡，才能避免把宋辽夏金元的多政权历史改写成单一中心的必然过程。

账本条目\texttt{YUAN-1237-REGNAL-YEAR}所记1237年的“蒙古以窝阔台在位第10年作为诏令、赋役与外交文书的纪年框架”，属于蒙古的历史语境。条目所示前因是新岁颁历、年号和纪年是中枢与地方行政沟通的共同前提，其可见后果为为同年政令、战事和地方材料提供日期锚点，不能单独推知具体政策执行。这一变化若进入区域生活，必须通过道路、文书、军队、市场、家庭和地方中介才会发生；政权易手或法令发布不等于所有城镇、乡村与边地在同一时间获得同样经验。

本条所列来源为《元史》〈本纪〉及相关志；《元朝秘史》；《元典章》或《通制条格》相关条。它们能够支持特定的年序、制度语言、政治行动或局部遗存，却不能无条件化为社会总量。账本特别提醒“纪年框架可靠；该记录仅确证政权纪年连续，年内具体事项须由本年条另证”。因此，军报的兵数、条约的名分、官府的族群分类、判牍中的道德判断以及后出的叙述都须回到产生它们的场景，而不可直接代表全域动机或人口。

将事件与地方层次连接时，应同时考察资源怎样被征集、信息怎样传递、迁徙怎样改变户籍和劳动、宗教与亲属网络怎样提供或限制协作。现有证据往往只照亮少数地点；保留这种不均衡，才能避免把宋辽夏金元的多政权历史改写成单一中心的必然过程。

账本条目\texttt{YUAN-1238-REGNAL-YEAR}所记1238年的“蒙古以窝阔台在位第11年作为诏令、赋役与外交文书的纪年框架”，属于蒙古的历史语境。条目所示前因是新岁颁历、年号和纪年是中枢与地方行政沟通的共同前提，其可见后果为为同年政令、战事和地方材料提供日期锚点，不能单独推知具体政策执行。这一变化若进入区域生活，必须通过道路、文书、军队、市场、家庭和地方中介才会发生；政权易手或法令发布不等于所有城镇、乡村与边地在同一时间获得同样经验。

本条所列来源为《元史》〈本纪〉及相关志；《元朝秘史》；《元典章》或《通制条格》相关条。它们能够支持特定的年序、制度语言、政治行动或局部遗存，却不能无条件化为社会总量。账本特别提醒“纪年框架可靠；该记录仅确证政权纪年连续，年内具体事项须由本年条另证”。因此，军报的兵数、条约的名分、官府的族群分类、判牍中的道德判断以及后出的叙述都须回到产生它们的场景，而不可直接代表全域动机或人口。

将事件与地方层次连接时，应同时考察资源怎样被征集、信息怎样传递、迁徙怎样改变户籍和劳动、宗教与亲属网络怎样提供或限制协作。现有证据往往只照亮少数地点；保留这种不均衡，才能避免把宋辽夏金元的多政权历史改写成单一中心的必然过程。

账本条目\texttt{YUAN-1239-REGNAL-YEAR}所记1239年的“蒙古以窝阔台在位第12年作为诏令、赋役与外交文书的纪年框架”，属于蒙古的历史语境。条目所示前因是新岁颁历、年号和纪年是中枢与地方行政沟通的共同前提，其可见后果为为同年政令、战事和地方材料提供日期锚点，不能单独推知具体政策执行。这一变化若进入区域生活，必须通过道路、文书、军队、市场、家庭和地方中介才会发生；政权易手或法令发布不等于所有城镇、乡村与边地在同一时间获得同样经验。

本条所列来源为《元史》〈本纪〉及相关志；《元朝秘史》；《元典章》或《通制条格》相关条。它们能够支持特定的年序、制度语言、政治行动或局部遗存，却不能无条件化为社会总量。账本特别提醒“纪年框架可靠；该记录仅确证政权纪年连续，年内具体事项须由本年条另证”。因此，军报的兵数、条约的名分、官府的族群分类、判牍中的道德判断以及后出的叙述都须回到产生它们的场景，而不可直接代表全域动机或人口。

将事件与地方层次连接时，应同时考察资源怎样被征集、信息怎样传递、迁徙怎样改变户籍和劳动、宗教与亲属网络怎样提供或限制协作。现有证据往往只照亮少数地点；保留这种不均衡，才能避免把宋辽夏金元的多政权历史改写成单一中心的必然过程。
